% GENERATED by scripts/manual_tables.py from the adapter registry. Do not edit:
% rerun the script. Hand-written text: blocks/catalogue/minirocket/
\ubblock{catalogue.minirocket}%
  {MiniRocket}% what the app calls it
  {MiniRocket shape fingerprint (WindowSet -> WindowSet + features)}% registry display name
  {encode}% insert-modal tab
  {WindowSet}% input
  {WindowSet}% output
  {Adapters/catalogue\_minirocket.py}% code
  {catalogue/minirocket}% hand-written folder under blocks/
\ubblockdescription{catalogue.minirocket}{Gives every window a shape fingerprint: about ten thousand tiny fixed pattern detectors, and for each the fraction of the window where it fires (MiniRocket, Dempster, Schmidt \& Webb 2021). A classifier can learn from it in seconds --- it answers 'can this label be learned from the raw shape at all?', but the numbers mean nothing one by one. Fitted on the training windows only (the split column), never on labels; each window z-normalised first. Without a split (Window matrix on the labels' 600/200 grid) it fits on every window and says so: acceptable there, because the fit only reads cut-offs from window shapes and never sees a label. Tune n\_features first: it sets the table's size and the time.}
\ubblockflags{catalogue.minirocket}{%
  \ubflag{derive}{}%
  \ubflag{estimate}{}%
}
\ubblocksideinputs{catalogue.minirocket}{%
}
\ubblockparams{catalogue.minirocket}{%
  \ubparam{n\_\allowbreak{}features}{10000}{whole number from 84 to 50000}{How many fingerprint numbers per window (the paper's 10,000; aeon rounds down to a multiple of 84, and the real count is shown). Tune this first: fewer is faster and smaller --- the table is windows \ensuremath{\times} features \ensuremath{\times} 4 bytes, so a whole channel on the 10-minute labels' grid (\textasciitilde{}13,000 windows) needs about 1,000 to stay under max\_cells. But fewer features lose real accuracy on weak shapes: a 'no' at 1,000 is not a 'no' --- the decisive run is 10,000 on a shorter span.}%
  \ubparam{max\_\allowbreak{}dilations}{32}{whole number from 1 to 128}{Most spacings per pattern detector (the paper's 32): how far apart the 9 taps of a detector may be spread, up to the whole window.}%
  \ubparam{seed}{0}{whole number, at least 0}{Which random training windows the detectors' cut-offs are read from; the same seed gives the same fingerprint.}%
  \ubparam{normalise}{yes}{yes or no}{z-normalise each window first (minus its mean, over its spread), so the fingerprint is about shape, not level or size --- as the paper's data are. Off: the window as it is.}%
  \ubparam{max\_\allowbreak{}cells}{50000000}{whole number, at least 1}{Refuse a table of more than this many numbers (windows \ensuremath{\times} features; 50 M = 200 MB): a memory guard. Lower n\_features before raising it.}%
  \ubparam{keep\_\allowbreak{}upstream\_\allowbreak{}features}{yes}{yes or no}{Keep the measures the windows already carried (Catch22, entropies) beside the fingerprint. Off: only the fingerprint and bookkeeping (split, pool and shape-file columns) go on, so a classifier learns from the shape fingerprint alone --- what an arm compared against Catch22 needs.}%
}
